@@M0@@, and if @@M1@@, @@M2@@. This does define a finite filtration on @@M3@@.

We define a homomorphism:
@@M4@@
by putting @@M5@@, which visibly does not depend on the choice of the @@M6@@ lifting the given sections @@M7@@ of @@M8@@. When the exact sequence is locally split, @@M9@@ is an isomorphism, since locally @@M10@@.

When @@M11@@ are locally direct summands of free Modules of finite type, the @@M12@@ are locally direct summands of free Modules of finite type, and therefore so are the @@M13@@; we may thus write in @@M14@@:
@@M15@@, giving additivity.

Let @@M16@@, @@M17@@ be two commutatively ringed topoi, and @@M18@@ a morphism of ringed topoi. We know (Exp. IV, 2.7) that @@M19@@ defines a ring morphism on the corresponding @@M20@@, namely @@M21@@. Since inverse image commutes with exterior powers, the corresponding morphism @@M22@@ on the @@M23@@ is a @@M24@@-homomorphism for the structure just defined.

Let us mention two special cases:
\begin{itemize}[label=-]
\item Let @@M25@@ be a scheme; @@M26@@ is thus endowed with a pre-@@M27@@-ring structure; this is the case that will be used in the rest of the seminar, and it may be regarded as a ``model'' for the general theory of
